\chapter{Conclusions}
\label{chap:conclusion}

Obtaining more markers does not necessarily mean that every
marker, or every available individual, contributes equally to a
classification problem. Livestock SNP datasets can contain tens of
thousands of correlated variables while the number of genotyped
animals remains comparatively small. The work, therefore, approaches
breed classification as a problem of learning, compression, and
interpretation at the same time.

The methodology connects these three levels. A leakage-safe genomic
pipeline defines which information is available during development;
variational and self-supervised contrastive models provide two ways of
constructing compact representations of complete genotype profiles;
classical population-genetics statistics, machine-learning feature
importance, and neural attribution methods, then move from the learned
representations back to individual markers.

The four questions share the same preprocessing and data-partitioning
rules, while the manipulated factor changes from one experiment to the
next: the number of training individuals for sample efficiency and the
number of markers for panel reduction.

The final evidence can be summarized in four points. RQ1 (breed
classification). Compact genomic representations classify livestock
breeds very accurately without species-specific tuning: Macro-F1
$0.9703$ and $0.9496$ on the caprine held-out test and $0.9555$ and
$0.9148$ on the independent ovine held-out test for the supervised and
self-supervised models respectively. RQ2 (sample efficiency). A
small training population suffices: in the caprine cohort
$N_{\min}=10$ individuals per breed retained $98.5\%$ of the reference
performance, and in the ovine cohort the same sample size retained
$98.6\%$. RQ3 (marker efficiency). The marker panel can be reduced by
two orders of magnitude under the frozen Random Forest evaluator: $100$
markers selected by a chi-square association test, or $200$ selected by
$F_{ST}$, match the full post-QC panel, while neural attributions
require larger panels ($500$ markers for SHAP and $2000$--$5000$ for
Integrated Gradients). RQ4 (marker-selection strategy). Neural
attribution is not a substitute for feature selection under that
evaluator: the
top-marker lists of the four families overlap little (mostly close to
chance for the neural methods), and their downstream utility differs
accordingly. The robustness analyses qualify this conclusion: on the caprine
cohort, when the panels are evaluated by a linear or boosting
classifier, or by retraining the explained model on the panel itself,
SHAP-selected panels become the most marker-efficient family, whereas
in the ovine cohort, the classical ordering is preserved under every
evaluator.
Cross-species validation supports the representation pipeline but
shows that its advantage over classical baselines is
dataset-dependent.

\section*{Limitations}

These conclusions are bounded by the available data and the
experimental protocol.

Breed classification is not equivalent to discovering causal
biological variants. A SNP can be useful for population discrimination
because it tags population structure or correlates with other loci,
without being a functional cause of breed identity.

Marker importance is also not unique in high-LD genomic data.
Several correlated markers may carry interchangeable predictive
information, so disagreement between rankings does not
automatically imply that one method is incorrect.

The number and distribution of genotyped individuals constrain what
can be concluded about sample efficiency. The operational threshold
identified here depends on the selected breeds, their genomic
diversity, marker density, preprocessing pipeline, learning
architecture, and retained-performance criterion; the resulting value
of $N_{\min}$ should not be interpreted as a universal biological
minimum number of animals required for breed classification.

The cross-species analysis is also methodological. ADAPTmap and
Sheep HapMap differ in species, marker space, and classification
problem; the second domain tests whether the behavior of the pipeline
can be reproduced in another livestock setting, and it does not assume
identical informative loci or directly comparable raw scores.

\section*{Future work}

Several extensions could strengthen the current study.

Biological annotation of the most stable and predictive SNPs is a
natural next step. Connecting computational rankings to genes,
regulatory regions, and known population-history signals would help
distinguish predictive markers from biologically interpretable
candidates.

Ranking uncertainty could be studied more explicitly.
Bootstrap procedures, repeated sampling, and methods that operate on
LD blocks rather than individual SNPs could provide a more realistic
view of feature importance when correlated markers are
interchangeable.

The representation-learning models could also be extended. The
variational formulation may be investigated with alternative
regularization strategies or multi-view genomic information, while
the contrastive model could be studied with augmentation policies that
are designed more explicitly around biological invariances.

Open-set breed assignment is a further extension: the model could be
asked to reject individuals that do not belong to any reference breed,
using the geometry of the learned sphere or the reconstruction error
as a rejection score.

Synthetic genomic data generation for data-scarce populations is
another possible extension. Generative models could, in principle, be used
to augment breeds represented by only a small number of real
individuals and to investigate whether synthetic genotype profiles can
recover part of the classification performance lost under severe
sample scarcity.

Such an extension would require more than generating numerically
plausible genotype vectors. Synthetic individuals should preserve
relevant population-genetic properties, including allele-frequency
distributions, linkage-disequilibrium structure, and within-breed
genomic diversity. Particular care would be required when very few
real animals are available, since a generative model trained on an
extremely limited population could reproduce close variants of the
observed individuals instead of learning a representative genomic
distribution.

The sample-efficiency analysis provides a basis for such future work. Once the range in which
classification performance deteriorates has been identified, synthetic
augmentation could be evaluated specifically in that data-scarce
regime by comparing models trained on real individuals alone with
models trained on the same real individuals supplemented by synthetic
genomic profiles.

The reduced-panel analysis could also be moved closer to an
applied traceability setting. Independent external cohorts, additional
breeds, and genotyping platforms with different marker densities would
provide stronger evidence about how well a computationally selected
panel transfers outside the datasets used during model development.

\section*{Practical impact}

A practical motivation is the possibility of replacing an
unnecessarily large genomic representation with a smaller one without
losing the information needed for a defined task. If reliable
reduced panels can be identified reproducibly, they may provide a
computational starting point for more economical breed-assignment or
traceability workflows.

Such a transition requires caution. A computational ranking is not by
itself a validated diagnostic panel, and independent biological and
technical validation would be required before deployment. The
contribution of the thesis is therefore best viewed as a
methodological bridge between high-density genomic data, modern
representation learning, and the design of compact, testable genomic
subsets.

The complete experimental pipeline that implements the protocol
described in this thesis, together with the frozen configurations and
the evaluation scripts, is available at
\url{https://github.com/lorenzosp01/genomic-dataset-reduction}.

\section*{AI usage statement}

Generative AI tools were used during the preparation of this thesis as
assistive instruments for software development and for the linguistic
and editorial revision of the manuscript. All scientific decisions,
the experimental design, the data analysis, and the interpretation of
the results were made by the author, who reviewed and verified every
AI-assisted output and takes full responsibility for the content of
the thesis.
